\documentclass[12pt]{article}
\usepackage[utf8]{inputenc}
\usepackage[T1]{fontenc}
\usepackage[margin=1in]{geometry}
\usepackage{setspace}
\usepackage{xcolor}
\usepackage[colorlinks=true, urlcolor=blue!60!black, linkcolor=blue!60!black]{hyperref}
\setlength{\parindent}{2em}
\setlength{\parskip}{1.25em}
\renewcommand{\baselinestretch}{1.0}
\begin{document}
\begin{center}
{\LARGE \textit{Green v. County School Board of New Kent County}, 391 U.S. 430 (1968)\par}
\vspace{0.8em}
{\large Case Brief --- Four-Page Edition\par}
\end{center}
\vspace{1em}

\section{Memory Jogger}

School boards that operated state-mandated dual school systems cannot satisfy their affirmative Fourteenth Amendment duty to eliminate segregation ``root and branch'' by adopting a ``freedom-of-choice'' plan that leaves racial identifiability intact when more effective alternatives such as geographic zoning are readily available (\href{https://supreme.justia.com/cases/federal/us/391/430/}{Justia}).

\section{Detailed Case Facts}

New Kent County, in rural eastern Virginia, had about 4,500 residents, half of them Black, and no residential segregation. Its roughly 1,300 pupils attended two combined elementary-and-high schools: New Kent School for white students and George W. Watkins School for Black students. There were no attendance zones; 21 buses ran overlapping routes countywide to carry children to their racially designated schools. Under Virginia's ``Massive Resistance'' Pupil Placement Act, children were automatically reassigned to their prior school, and by 1964 no student of either race had sought a transfer. After Black parents led by Charles C. Green sued in March 1965, the Board, to keep federal funds under Title VI, adopted a ``freedom-of-choice'' plan letting students choose annually, with nonchoosers reassigned to their prior school. In three years no white student chose Watkins, and 85\% of Black students remained there.

\section{Procedural History}

The Eastern District of Virginia approved the amended freedom-of-choice plan in 1966. The en banc Fourth Circuit, relying on its companion decision in \textit{Bowman v. County School Board of Charles City County}, affirmed approval of freedom of choice but remanded for a firmer faculty-desegregation order; Judges Sobeloff and Winter disagreed with upholding freedom of choice. The Supreme Court granted certiorari, heard argument April 3, 1968, with \textit{Raney} and \textit{Monroe}, and vacated and remanded on May 27, 1968.

\section{Judicial Votes}

Vacated and remanded, \textbf{9--0}.
\begin{itemize}
\item \textbf{Majority (9):} Justice Brennan, joined by Chief Justice Warren and Justices Black, Douglas, Harlan, Stewart, White, Fortas, and Marshall.
\item \textbf{Concurrences and dissents:} None.
\end{itemize}

\section{Holding}

A board that ran a state-compelled dual system at the time of \textit{Brown I} has an affirmative duty to convert promptly to a unitary, nonracial system in which discrimination is eliminated ``root and branch'' across student assignment, faculty, staff, transportation, extracurricular activities, and facilities (the \textit{Green} factors). Its plan must promise ``realistically to work now.'' Freedom of choice is not unconstitutional per se, but it is unacceptable when it fails to dismantle the dual system and readily available alternatives, such as zoning or grade consolidation, promise faster conversion. The district court must retain jurisdiction until state-imposed segregation is removed. The Court's suggestions about where choice might work and about specific remedies were dicta.

\section{Analysis of Opinions}

\textbf{Majority (Brennan).} New Kent ran the archetypal ``dual system,'' racially identifiable in all six operational dimensions. The Board argued that a race-neutral choice plan fully discharged its duty and that more would be ``compulsory integration.'' Brennan answered that opening the doors in 1965 ``merely begins, not ends, our inquiry,'' because \textit{Brown II} required eradicating segregation's vestiges. The Board's eleven-year delay compounded the harm, and quoting \textit{Griffin}, ``[t]he time for mere `deliberate speed' has run out.'' Adopting Judge Sobeloff's view, Brennan treated freedom of choice as a means, not an end. In a residentially integrated county where zoning or splitting the schools by grade would work, a plan leaving 85\% of Black children at Watkins improperly shifted the burden of desegregation from the Board to Black families.

\textbf{Separate opinions.} None.

\section{Examples: Future Applications}

\begin{enumerate}
\item \textbf{Same-Side: Free-Transfer Escape Hatch.} A formerly segregated district draws integrating zones but lets any student transfer freely, and white students zoned into formerly Black schools transfer out. Like New Kent's plan, the option perpetuates the dual system (as in the companion \textit{Monroe} case).
\item \textbf{Same-Side: Residual Segregation in Faculty and Facilities.} A district balances enrollment but keeps 95\% of Black teachers at the formerly Black campus, with inferior athletic facilities and segregated clubs. \textit{Green} requires eliminating racial identifiability across all six factors, so the district is not unitary.
\item \textbf{Opposite-Side: Purely De Facto Imbalance.} A suburban district in the North that never practiced intentional segregation develops racially imbalanced neighborhood schools through housing patterns. \textit{Green}'s duty is triggered only by prior de jure segregation.
\item \textbf{Opposite-Side: Resegregation After Unitary Status.} A district complies in good faith for twenty years, is declared unitary, and later adopts race-neutral neighborhood schools that become predominantly single-race. The remedial duty has ended absent new intentional discrimination.
\item \textbf{Fence-Sitter: Magnet-School Choice.} A formerly segregated county places magnet academies at formerly Black schools with free transportation; white enrollment there rises from 0\% to 30\% in a 50/50 district, though mandatory reassignment could achieve balance at the cost of white flight. A court must decide whether the plan ``offers real promise'' or whether a faster remedy makes it inadequate.
\end{enumerate}

\section{Critique}

\textbf{Originalist and textualist critiques.} Conservative scholars argue that \textit{Green} shifted, without textual warrant, from \textit{Brown I}'s colorblind ban on racial classification to a mandate of race-conscious ``compulsory integration.'' Measuring compliance by enrollment ratios invited decades of judicial management of local schools, busing, and racial balancing beyond Article III remedial limits.

\textbf{Progressive and doctrinal perspectives.} Civil rights advocates praise \textit{Green} for finally giving \textit{Brown} teeth by demanding results rather than paper neutrality. Yet tying the duty to prior de jure systems let later decisions confine it, shield suburbs from interdistrict remedies, and dissolve decrees on findings of unitary status, contributing to resegregation. The case embodies the clash between anticlassification and antisubordination.

\section{Key Quotations}

\begin{quote}
``School boards such as the respondent then operating state-compelled dual systems were nevertheless clearly charged with the affirmative duty to take whatever steps might be necessary to convert to a unitary system in which racial discrimination would be eliminated root and branch.''
\end{quote}
\noindent --- Justice Brennan, for a unanimous Court, 391 U.S. at 437--38.

\begin{quote}
``The burden on a school board today is to come forward with a plan that promises realistically to work, and promises realistically to work now.''
\end{quote}
\noindent --- Justice Brennan, 391 U.S. at 439.

\begin{quote}
``Rather than further the dismantling of the dual system, the plan has operated simply to burden children and their parents with a responsibility which \textit{Brown II} placed squarely on the School Board. The Board must be required to formulate a new plan and, in light of other courses which appear open to the Board, such as zoning, fashion steps which promise realistically to convert promptly to a system without a `white' school and a `Negro' school, but just schools.''
\end{quote}
\noindent --- Justice Brennan, 391 U.S. at 441--42.

\end{document}
